\chapter{Underactuation and Passive Dynamics}
\label{ch:underactuation}

\begin{laymansbox}{How One Input Moves Several Links}{working_with_physics}
Moving the hands can accelerate a club even when a simplified model applies
no wrist torque. The club's response depends on its orientation, existing
velocity, gravity, and the acceleration of the hands. This is a useful way to
study release. It does not establish that a golfer's wrists are freely hinged,
that the arms must actively brake, or that the resulting motion is stable or
optimal. Those are separate questions that the model and measurements must answer.
\end{laymansbox}

\section{Instantaneous Authority and Motion Over Time}

In independent configuration coordinates $q\in\mathbb R^n$, write
\begin{equation}
M(q)\ddot q+h(q,\dot q)=B(q)u,\qquad M=M^T\succ0.
\end{equation}
Here $h=C(q,\dot q)\dot q+\nabla V(q)$ for an ideal rigid mechanical model;
damping, elastic forces, and other loads must be added when retained.
At fixed state the attainable acceleration set is
\begin{equation}
\mathcal A(q,\dot q)=-M^{-1}h+M^{-1}B\mathcal U,
\end{equation}
where $\mathcal U$ is the admissible input set. With unrestricted inputs its
affine dimension is $r=\operatorname{rank}B$. Underactuation means $r<n$.
When $m<n$, an $n\times m$ matrix can have full column rank $m$ while lacking
full row rank. Multiple muscles do not guarantee independent joint torques:
moment arms, force signs, activation, and force limits constrain the attainable set.

Instantaneous authority is different from reachability over time. Coupled
dynamics can move an unactuated coordinate. A full-row-rank model also cannot
follow every trajectory when contact, torque, power, speed, or state constraints
are imposed. A walking foot has no moment at an ideal point contact, whereas
distributed foot contact can support a bounded moment. The contact model matters.

\section{A Coordinate-Consistent Double Pendulum}

Let $\theta_1$ and $\theta_2$ be absolute angles of the arm and club measured
from downward vertical in one inertial plane. Let $\tau_s$ be shoulder torque
and $\tau_w$ wrist torque exerted on the club. Virtual work gives
\begin{equation}
\delta W=\tau_s\delta\theta_1+\tau_w(\delta\theta_2-\delta\theta_1),
\qquad Q_\theta=(\tau_s-\tau_w,\tau_w)^T.
\end{equation}
The corresponding actuator power is
$\tau_s\dot\theta_1+\tau_w(\dot\theta_2-\dot\theta_1)$.
For relative coordinates $q=(\theta_1,\theta_2-\theta_1)^T$, let
\begin{equation}
\theta=Tq,\qquad T=\begin{pmatrix}1&0\\1&1\end{pmatrix}.
\end{equation}
Then $Q_q=T^TQ_\theta=(\tau_s,\tau_w)^T$ and $M_q=T^TM_\theta T$.
For this constant transformation, $h_q=T^Th_\theta$. The shoulder-only map
$(1,0)^T$ is valid in both conventions when $\tau_w=0$; the mass and bias terms
must still use the chosen coordinates. Wrist power uses relative angular speed.

For a transparent worked model, use massless links of lengths $l_1,l_2$,
point masses $m_1$ at the wrist and $m_2$ at the club tip, and positive gravity
magnitude $g$. Set $\Delta=\theta_1-\theta_2$ and $b=m_2l_1l_2$. In absolute angles,
\begin{equation}
M_\theta=\begin{pmatrix}
(m_1+m_2)l_1^2&b\cos\Delta\\
b\cos\Delta&m_2l_2^2
\end{pmatrix},
\end{equation}
\begin{equation}
h_\theta=\begin{pmatrix}
b\sin\Delta\,\dot\theta_2^2+(m_1+m_2)gl_1\sin\theta_1\\
-b\sin\Delta\,\dot\theta_1^2+m_2gl_2\sin\theta_2
\end{pmatrix}.
\end{equation}
These equations follow from $L=\tfrac12\dot\theta^TM_\theta\dot\theta-V$ with
$V=-(m_1+m_2)gl_1\cos\theta_1-m_2gl_2\cos\theta_2$.
They omit distributed limb inertia, a moving shoulder, shaft bending, and
three-dimensional grip mechanics. Their purpose is an inspectable mechanism.

\subsection{What Arm Deceleration Does and Does Not Imply}

With zero wrist torque, the second row is
\begin{equation}
M_{21}\ddot\theta_1+M_{22}\ddot\theta_2+h_2=0,
\qquad
\ddot\theta_2=-\frac{M_{21}\ddot\theta_1+h_2}{M_{22}}.
\label{eq:pend_unactuated}
\end{equation}
The positive denominator does not determine the sign of this expression.
$M_{21}=b\cos\Delta$ changes sign with geometry, and the velocity and gravity
terms remain. With $M_{22}=1$, $h_2=0$, and $\ddot\theta_1=-2$,
$M_{21}=0.4$ gives $\ddot\theta_2=0.8$ whereas $M_{21}=-0.4$ gives $-0.8$.
Both matrices with $M_{11}=2$ are positive definite. These are a dimensionally
consistent rotational-inertia example, not measured golfer parameters.

Arm deceleration can therefore accompany club acceleration in suitable states.
It is not a necessary instruction to actively brake the arms: reaction torques
can decelerate an arm while shoulder torque remains positive. Angular momentum
of the arm-plus-club system also changes under external shoulder torque and
gravity. Calling this an optimal momentum transfer requires an objective,
admissible inputs, and an optimization result that these equations alone do not supply.

For the ideal conservative model,
\begin{equation}
\frac{d}{dt}(T_{\rm kinetic}+V)
=\tau_s\dot\theta_1+\tau_w(\dot\theta_2-\dot\theta_1).
\end{equation}
Zero wrist actuator power does not mean zero energy flow into the club.
Forces at the moving wrist can do work on the club while doing opposite work
on the arm. Whole-system and segment energy budgets answer different questions;
include translational boundary power and avoid counting internal power twice.

\section{The Annihilator Test for Feasible Trajectories}

On a constant-rank coordinate patch, let $N_B(q)\in\mathbb R^{(n-r)\times n}$
have rows spanning the left nullspace of $B$. Thus $N_BB=0$. A twice
differentiable configuration path with its stated timing requires the load
\begin{equation}
b_*(t)=M(q_*)\ddot q_*+h(q_*,\dot q_*).
\end{equation}
The condition
\begin{equation}
N_B(q_*)b_*=0
\label{eq:unactuated_dynamics}
\end{equation}
is equivalent to $b_*\in\operatorname{range}B$: some unrestricted algebraic
input produces that load. One solution is $u_*=B^+b_*$, with possible additions
from $\ker B$. The required input must also belong to $\mathcal U$.
For example, $B=(1,0)^T$ and $b_*=(10,0)^T$ pass the annihilator test but fail
an input bound $|u|\leq1$. Input-rate and actuator-state constraints require
additional dynamical checks. A locally smooth constant-rank pseudoinverse does
not remove these restrictions.

This is a constraint on position, velocity, and acceleration along a trajectory,
not generally a fixed subspace of state space. Changing the timing changes
$\dot q_*$ and $\ddot q_*$, so the same geometric path may pass or fail.
Contact reactions must first be eliminated consistently or retained as additional
unknowns with their friction and unilateral-contact conditions.

\section{Partial Feedback Linearization and Zero Dynamics}

For an ideal collocated partition with $B=(I,0)^T$, write
\begin{align}
M_{aa}\ddot q_a+M_{au}\ddot q_u+h_a&=u,\\
M_{ua}\ddot q_a+M_{uu}\ddot q_u+h_u&=0.
\end{align}
Choosing the acceleration command $v_a$ requires
\begin{equation}
u=(M_{aa}-M_{au}M_{uu}^{-1}M_{ua})v_a
  +h_a-M_{au}M_{uu}^{-1}h_u.
\end{equation}
The Schur complement is positive definite when $M\succ0$. Exact cancellation
assumes an accurate model and an admissible computed input. Saturation and
uncertain loads invalidate the ideal equality $\ddot q_a=v_a$.

For output $e_a=q_a-q_a^*(t)$, zero dynamics are the remaining dynamics on
$e_a=\dot e_a=0$ under the input that keeps this set invariant. They depend on
the output and its reference, and are generally time varying for tracking.
They are not identical to the motion obtained by switching all actuators off.

No phase of a generic swing is automatically attracting. Integrate a specified
nominal solution and the variational equation
\begin{equation}
\delta\dot z=A_z(t)\delta z,\qquad
\dot\Phi_z=A_z(t)\Phi_z,\quad\Phi_z(t_0)=I.
\end{equation}
Specify the output, comparison time or event, and the metric used to scale
angle and velocity errors. Singular values of the appropriately scaled
transition map measure finite-time amplification. A time-varying positive
metric $W(t)$ satisfying
\begin{equation}
\dot W+A_z^TW+WA_z\preceq-2\lambda W
\end{equation}
certifies differential decay on the interval where it holds, with metric bounds
needed to translate it into a coordinate norm. Frozen-time eigenvalues and
large club speed alone do not supply this certificate. A verified invariant
funnel additionally checks its moving boundary, input limits, and disturbances.

\section{Accessibility Is Not Forward-Time Reachability}

For driftless systems $\dot x=\sum_i g_i(x)u_i$ with reversible inputs, a
bracket-generating distribution on a connected manifold gives the horizontal
connectivity described by Chow--Rashevskii. This does not immediately extend
to a drifting mechanical system $\dot x=f(x)+\sum_i g_i(x)u_i$ with bounded
inputs and a fixed downswing duration.

An orbit generated by vector-field flows permits both positive and negative
flow times. Forward reachable sets do not reverse an unavoidable drift.
For $\dot x=1$, $\dot y=u$, $|u|\leq1$, the fields span the plane, yet from
the origin $x(T)=T$ and $|y(T)|\leq T$. No negative $x$ is reachable for
$T\geq0$. Lie-algebra rank can support local accessibility under appropriate
regularity and input assumptions; it does not certify global or small-time
local controllability. Smooth orbit theorems also require care beyond a
pointwise Lie-bracket rank test.

Abnormal extremals have a zero cost multiplier in Pontryagin's principle;
singular arcs require separate analysis of the Hamiltonian's input dependence.
Neither is diagnosed by visually rapid wrist release. Establish them for a
specified optimal-control problem before assigning them a physiological meaning.

\section{Connecting the Reduced Model to a Real Swing}

The useful hypothesis is that proximal motion and earlier work prepare a state
from which mechanical coupling contributes to club delivery. Testing it requires
measured wrist and hand wrenches, consistent angular coordinates, a moving-base
model, and uncertainty estimates. A passive hinge is one comparison model.
Compare it with bounded wrist torques and compliant grips on the same task;
then assess speed, orientation, dispersion, and effort separately.

Shaft deformation stores energy and changes the contact-point velocity. Muscle
activation affects force capacity and impedance. Feedback changes responses to
perturbations, while preparation changes the nominal trajectory itself. These
mechanisms can coexist. A zero-torque comparison does not identify their unique
biological contributions, and a simulated advantage must survive measurement
and model-sensitivity checks before becoming a coaching recommendation.

\section{Exercises}
\enlargethispage{3\baselineskip}

\begin{enumerate}
\item Verify the mass and bias terms from the Lagrangian and check the power
identity with nonzero wrist torque in both coordinate systems.
\item Reproduce the coupling-sign counterexample. Include gravity and nonzero
angular rates and identify when club acceleration changes sign.
\item Construct $N_B$ for a rank-deficient actuator map. Give a load that passes
the range test but violates an actuator bound.
\item For $\ddot q_u+cq_u=\sin t$, subtract a particular solution. Classify
homogeneous perturbations for $c>0$, $c=0$, and $c<0$; examine resonance at $c=1$.
Explain why bounded perturbations do not guarantee bounded forced motion.
\item Specify an experiment that distinguishes passive-hinge prediction error
from measurement noise, wrist actuation, and shaft compliance.
\end{enumerate}
